\section{FORMAL RESTATEMENT}
\textbf{Erd\H{o}s \#286; an exact-length deduction and a finite exception, 2026-09-24.}
The displayed question asks for $k$ distinct positive integers in an interval of width $(e-1+o(1))k$, whose reciprocals sum to $1$. Its meaningful asymptotic assertion is for every sufficiently large integer $k$. The phrase $k\geq2$ has an exception at $k=2$, as proved below.

\section{QUICK LITERATURE AND CONTEXT CHECK}
The problem page attributes the affirmative asymptotic result to Croot. Martin's exact-length theorem gives a convenient, stronger upper bound on the largest denominator and therefore also proves the interval assertion as literally formulated, since the containing interval may be wider than the span of the selected denominators. This deduction uses Martin's theorem as an external input; it does not purport to reproduce Croot's short-interval construction.

\section{OBJECTIVE AND ATTACK PLAN}
Keep the number of terms equal to the prescribed $k$. Place the resulting denominators in an explicitly chosen interval and check that its width is exactly the requested leading constant. Separate the impossible two-term case from the large-$k$ claim.

\section{WORK}
\textbf{The finite exception.} If $1/u+1/v=1$, then $(u-1)(v-1)=1$. Positive denominators therefore force $u=v=2$, which violates distinctness. Thus there is no two-term example, regardless of interval width. Every length $k\geq3$ is feasible: start with $1=1/2+1/3+1/6$ and repeatedly replace the term with largest denominator $m$ by $1/(m+1)+1/[m(m+1)]$. The new denominators exceed all retained ones.

\textbf{Declared exact-length input.} Martin's Proposition 5, with target rational $r=1$, provides for every sufficiently large integer $k$ a set $E_k$ of exactly $k$ distinct positive integers satisfying
\[
\sum_{n\in E_k}\frac1n=1,\qquad
\max E_k\leq\frac e{e-1}k
+O\!\left(\frac{k\log\log k}{\log k}\right).
\tag{1}
\]
Its denominator smoothness hypothesis is $P^*(1)=1<k/(\log k)^{22}$ and hence is satisfied eventually. The proof of that construction is the external dependency.

There is a strict gap between the two constants:
\[
\frac e{e-1}<e-1
\quad\Longleftrightarrow\quad e^2-3e+1>0.
\tag{2}
\]
For example $e>1+1+1/2+1/6=8/3$, the polynomial $x^2-3x+1$ is increasing for $x>3/2$, and its value at $8/3$ is $1/9>0$. If $\eta=(e-1)-e/(e-1)>0$, the error in (1) is at most $\eta k/2$ for sufficiently large $k$. Therefore
\[
\max E_k<(e-1)k.
\]
Take
\[
I_k=[1,\,1+(e-1)k].
\tag{3}
\]
This interval contains every denominator in $E_k$ and has width exactly $(e-1)k$, which is of the required form $(e-1+o(1))k$. If integer endpoints are desired, use $[1,\lceil1+(e-1)k\rceil]$, whose width differs by less than one.

For the finitely many feasible lengths before (1) applies, any containing interval may be enlarged if necessary. Arbitrary changes at finitely many $k$ do not change the asymptotic width. They cannot repair $k=2$, where no representation exists.

\section{VERIFICATION AND SCOPE}
The variable $k$ in (1) is the number of summands, so no unjustified conversion from a lower-denominator parameter is involved. The assertion concerns a containing interval, not a requirement that its endpoints equal the least and greatest chosen denominators. The intended asymptotic statement is proved; the stronger literal assertion including two distinct terms is false. No novelty is claimed.

\section{SOURCES}
\begin{itemize}
\item Current statement and Croot attribution: \url{https://www.erdosproblems.com/286}.
\item G. Martin, \emph{Denser Egyptian Fractions}, Proposition 5 and Theorem 2: \url{https://arxiv.org/pdf/math/9811112}.
\item E. S. Croot III, \emph{On unit fractions with denominators in short intervals}: \url{https://arxiv.org/abs/math/9904181}.
\end{itemize}
\section{CONCLUSION}
\textbf{Attempt status: solved, with the $k=2$ exception stated explicitly.}
The large-$k$ interval assertion follows from Martin's exact-length theorem.
\textbf{Completion rate estimate: 100\% of the asymptotic assertion, relative to that theorem.}
